\title{When Is the Twisted Signing of an Even Cycle Square Spectrally Optimal?}
\ifdefined\TargetAAnonymous
  \author{Anonymous}
\else
  \author{[AUTHOR NAME]}
\fi
\date{}
\maketitle

\ifdefined\TargetAAnonymous\else
\begin{center}
\small
\mbox{[AFFILIATION]}\\
\mbox{[DEPARTMENT]}, \mbox{[INSTITUTION]}\\
\mbox{[CITY]}, \mbox{[COUNTRY]}\\[3pt]
Corresponding author: [CORRESPONDING AUTHOR]\\
Email: [EMAIL]\quad ORCID: [ORCID]
\end{center}
\fi

\begin{abstract}
For even \(n\geq8\), let \(m_n\) be the minimum adjacency spectral radius
over all edge signings of the square \(C_n^2=C_n(1,2)\). A distinguished
twisted signing attains \(\rho_-(n)\), where
\(\rho_-(n)^2=4+2\cos(2\pi/n)+2\cos(4\pi/n)\), at every even order.  We prove
that the conjectured equality \(m_n=\rho_-(n)\) holds exactly for
\(n\in\{8,10,12,14,16,18,20,22,24,26,28,30,34,36,38,42,44,46\}\); equivalently,
failure occurs at \(32\), at \(40\), and at every even \(n\geq48\).
The structural mechanism is a period-eight reference phase and its six-gap
phase slip \(G_6\). The reference squared spectral edge is \(\eta<8\), while
\(G_6\) supports an isolated top squared eigenvalue \(c_6\in(\eta,8)\) of
multiplicity two and uniquely minimizes the edge among abnormal positive
single gaps. Charge-compatible separated copies, combined with discrete IMS
localization, prove eventual failure; finite exact verification identifies
the sharp onset at \(48\). The supplement records the certificates, their
verification contracts, and which objects are independently reconstructed.
\end{abstract}

\noindent\textbf{Keywords:} signed graph; spectral radius; cycle square;
switching equivalence; phase slip; computer-assisted proof

\medskip
\noindent\textbf{MSC 2020:} 05C50 (primary); 05C22, 68V05 (secondary)
